\chapter{Análisis del Circuito}
\section{Polarización del Diodo}
En esta actividad se realizará el análisis de polarización de un diodo utilizando un circuito compuesto por una fuente de tensión variable, una resistencia
 y un diodo conectados en serie. El objetivo es determinar el comportamiento del diodo para distintos valores de tensión de entrada aplicando la Ley de
 Tensiones de Kirchhoff (LTK).

A partir del circuito propuesto se calcularán la tensión en el diodo ($V_D$), la corriente que circula por él ($I_D$) y la potencia disipada ($P_D$).
 El análisis se efectuará considerando los distintos modelos estudiados para el diodo: ideal, simplificado y lineal.

 \begin{figure}[H]
    \centering
    \includegraphics[width=0.6\textwidth]{img/Analisis.png}
    \caption{Circuito de polarización del diodo.}
    \label{fig:diodo}
 \end{figure}

\subsection{Modelo ideal}
En el modelo ideal, el diodo se comporta como un cortocircuito cuando está polarizado en directa y como un circuito abierto cuando está polarizado en inversa.

\begin{itemize}
    \item En directa:$V_D=0V$ 
    \item En inversa: circuito abierto
\end{itemize}

Entonces, la corriente del diodo queda:
\[
I_D = \frac{V_1}{1000}
\]
\begin{table}[H]
\centering
\caption{Cálculos para Diodo de Silicio 1N4007.}
\begin{tabular}{c c c c}
\hline
$V_1$ [V] & $V_D$ [V] & $I_D$ [mA] & $P_D$ [mW] \\
\hline
0,5 &0 & 0,5&0 \\
1,0 &0 & 1&0 \\
3,0 &0 & 3&0 \\
5,0 &0 & 5&0 \\
\hline
\end{tabular}
\end{table}

\vspace{0.5cm}

\begin{table}[H]
\centering
\caption{Cálculos para Diodo de Germanio 1N60.}
\begin{tabular}{c c c c}
\hline
$V_1$ [V] & $V_D$ [V] & $I_D$ [mA] & $P_D$ [mW] \\
\hline
0,2 &0 &0,2 &0 \\
0,5 &0 &0,5 &0 \\
1,0 &0 &1 &0 \\
3,0 &0 &3 &0 \\
\hline
\end{tabular}
\end{table}

\subsection{Modelo simplificado}

En el modelo simplificado se considera una caída de tensión fija en el diodo cuando conduce.

Para el diodo de Silicio 1N4007 se considera $V_D \approx 0.7V$, mientras que para el diodo de Germanio 1N60 se utiliza $V_D \approx 0.3V$.
Entonces, la corriente del diodo queda:

\[
I_D = \frac{V_1 - V_D}{1000}
\]

Siempre que:$V_1 > V_D$.
En caso contrario:$I_D = 0$(Se lo considera como una llave abierta.)
\begin{table}[H]
\centering
\caption{Cálculos para Diodo de Silicio 1N4007.}
\begin{tabular}{c c c c}
\hline
$V_1$ [V] & $V_D$ [V] & $I_D$ [mA] & $P_D$ [mW] \\
\hline
0,5 &0,5 &0&0 \\
1,0 &0,7 &0,3 &0,21 \\
3,0 &0,7 &2,3&1,61 \\
5,0 &0,7 &4,3&3,01 \\
\hline
\end{tabular}
\end{table}

\vspace{0.5cm}

\begin{table}[H]
\centering
\caption{Cálculos para Diodo de Germanio 1N60.}
\begin{tabular}{c c c c}
\hline
$V_1$ [V] & $V_D$ [V] & $I_D$ [mA] & $P_D$ [mW] \\
\hline
0,2 &0,2 &0 &0 \\
0,5 &0,3 &0,2 &0,06 \\
1,0 &0,3 &0,7 &0,21 \\
3,0 &0,3 &2,7 &0,81 \\
\hline
\end{tabular}
\end{table}

\subsection{Modelo lineal}

En el modelo lineal, además de la caída de tensión fija, se considera una resistencia interna del diodo:

\[
V_D = V_{\gamma} + I_D r_d
\]

Para el diodo de Silicio se toma generalmente $V_{\gamma}=0.7V$ y $r_d \approx 10\Omega$, mientras que para el diodo de Germanio se utiliza $V_{\gamma}=0.3V$ y $r_d \approx 5\Omega$.

Entonces, la corriente del diodo queda:

\[
I_D = \frac{V_1 - V_{\gamma}}{R_1 + r_d}
\]
\begin{table}[H]
\centering
\caption{Cálculos para Diodo de Silicio 1N4007.}
\begin{tabular}{c c c c}
\hline
$V_1$ [V] & $V_D$ [V] & $I_D$ [mA] & $P_D$ [mW] \\
\hline
0,5 &0,5 &0&0 \\
1,0 &0,7 &0,29 &0,20 \\
3,0 &0,7 &2,27&1,59 \\
5,0 &0,7 &4,25&2,98 \\
\hline
\end{tabular}
\end{table}

\vspace{0.5cm}

\begin{table}[H]
\centering
\caption{Cálculos para Diodo de Germanio 1N60.}
\begin{tabular}{c c c c}
\hline
$V_1$ [V] & $V_D$ [V] & $I_D$ [mA] & $P_D$ [mW] \\
\hline
0,2 &0,2 &0 &0 \\
0,5 &0,3 &0,199 &0,1 \\
1,0 &0,3 &0,697 &0,209 \\
3,0 &0,3 &2,687 &0,806 \\
\hline
\end{tabular}
\end{table}


\section{Circuitos Recortadores con Zener}

En esta sección se realizará el análisis de circuitos recortadores utilizando diodos zener. Se estudiará el comportamiento de la tensión de salida durante los semiciclos positivo y negativo de la señal de entrada, 
identificando qué diodo conduce en cada caso. Además, se graficarán las formas de onda de entrada y salida para observar el efecto de recorte producido por los diodos zener.

\begin{figure}[H]
    \centering
    \includegraphics[width=0.7\textwidth]{img/Analisis2.png}
    \caption{Circuitos recortadores con diodos zener.}
    \label{fig:zener}
\end{figure}

\subsection{Circuito a)}
Durante el semiciclo positivo, el diodo zener de $5.1V$ queda polarizado en directa, comportándose como un diodo convencional con una caída aproximada de $0.7V$. Por otro lado, el diodo zener de $20V$ queda polarizado en inversa.

Por lo tanto, la tensión mínima necesaria para que el circuito entre en la región de recorte es:

\[
V_{out} \approx 20V + 0.7V\approx 20.7
\]

Si    $V_{in} < 20.7V$
el diodo zener de $20V$ todavía no alcanza la región zener, por lo que se comporta como un diodo polarizado en inversa y la corriente 
es aproximadamente nula ($I_D \approx 0A$).


Aplicando LTK:

\[
V_{in} + I_D R - V_{out} = 0\]
\[V_{out} = V_{in}\]

Cuando $V_{in} > 20.7V$el diodo zener entra en ruptura y comienza el recorte de la señal.

Durante el semiciclo negativo, el diodo zener de $20V$ queda polarizado en directa, mientras que el diodo zener de $5.1V$ queda polarizado en inversa.

La tensión de recorte resulta:
\[
V_{sal} \approx -(5.1V + 0.7V)
\]
\[
V_{sal} \approx -5.8V
\]
Si $V_{in} > -5.8V$
los diodos no entran en conducción y:
\[
V_{out} = V_{in}
\]
Cuando $V_{in} < -5.8V$
el diodo zener de $5.1V$ entra en ruptura y comienza el recorte de la señal.

\begin{figure}[H]
    \centering
    \includegraphics[width=0.7\textwidth]{img/forma onda a.png}
    \caption{Forma de onda Circuito a}
    \label{fig:ZenerA}
\end{figure}



\subsection{Circuito b)}

Durante el semiciclo positivo, el diodo zener de $12V$ queda polarizado en directa, comportándose como un diodo convencional con una caída aproximada de $0.7V$. Por otro lado, el diodo zener de $6.8V$ queda polarizado en inversa.

Por lo tanto, la tensión mínima necesaria para que el circuito entre en la región de recorte es:

\[
V_{out} \approx 6.8V + 0.7V
\]
\[
V_{out} \approx 7.5V
\]

Si:$V_{in} < 7.5V$el diodo zener de $6.8V$ todavía no alcanza la región zener, por lo que se comporta como un diodo polarizado en inversa y$I_D \approx 0A$

Aplicando LTK:
\[
V_{in} + I_D R - V_{out} = 0
\]
\[
V_{out} = V_{in}
\]

Cuando $V_{in} > 7.5V$ el diodo zener de $6.8V$ entra en ruptura y comienza el recorte de la señal.
Por lo tanto, la salida queda limitada aproximadamente a:

\[
\boxed{V_{out} \approx 7.5V}
\]
Durante el semiciclo negativo, el diodo zener de $6.8V$ queda polarizado en directa, comportándose como un diodo convencional con una caída aproximada de $0.7V$. Por otro lado, el diodo zener de $12V$ queda polarizado en inversa.

La tensión de recorte resulta:

\[
V_{sal} \approx -(12V + 0.7V)
\]
\[
V_{sal} \approx -12.7V
\]

Si$V_{in} > -12.7V$
los diodos no entran en conducción y:
\[
V_{out} = V_{in}
\]

Cuando $V_{in} < -12.7V$ el diodo zener de $12V$ entra en ruptura y comienza el recorte de la señal.
Por lo tanto, la salida queda limitada aproximadamente a:

\[
\boxed{V_{out} \approx -12.7V}
\]

\begin{figure}[H]
    \centering
    \includegraphics[width=0.7\textwidth]{img/formaOndaB.png}
    \caption{Forma de onda Circuito b}
    \label{fig:ZenerB}
\end{figure}

\section{Análisis sobre parámetros de hoja de datos}
\subsection{Características Eléctricas}

\subsubsection*{Regímenes máximos}

\begin{itemize}

    \item $V_{RRM}$: Tensión inversa repetitiva máxima.  
    Es la máxima tensión inversa que el diodo puede soportar repetidamente sin dañarse.  
    \item $V_{RSM}$: Tensión inversa no repetitiva máxima.  
    Es la máxima tensión inversa instantánea que el diodo soporta durante transitorios.  
    \item $I_{FRMS}$: Corriente directa eficaz máxima.  
    Es la máxima corriente RMS permitida en conducción directa. 
    \item $I_{FSM}$: Corriente de sobrecarga máxima.  
    Corriente pico no repetitiva soportada durante un corto intervalo de tiempo. 

\end{itemize}


\subsubsection*{Regímenes característicos}

\begin{itemize}

    \item $V_{BR}$: Tensión de ruptura.
    Tensión inversa donde  el diodo entra en ruptura.  
    \item $I_R$: Corriente inversa.  
    Pequeña corriente que circula en polarización inversa.  
    \item $V_F$: Tensión directa.  
    Caída de tensión del diodo cuando conduce en directa. 
    \item $I_{RM}$: Corriente inversa máxima.  
    Máxima corriente inversa permitida. 
    \item $I_F$: Corriente directa.  
    Corriente que circula cuando está polarizado en directa.  

\end{itemize}


\subsection{Características Térmicas}

\begin{itemize}

    \item $T_J$: Temperatura de juntura.  
    Temperatura máxima permitida en la unión interna del semiconductor.  
    \item $T_{STG}$: Temperatura de almacenamiento.  
    Rango de temperaturas permitido para almacenar el dispositivo.  
    \item $T_A$: Temperatura ambiente.  
    Temperatura del entorno donde opera el dispositivo.  
    \item $R_{\theta JC}$: Resistencia térmica juntura-carcasa.  
    Indica la dificultad para disipar calor desde la juntura hacia la carcasa.  

\end{itemize}

\subsection{Características de uso}

\begin{itemize}

    \item $V_{RMS}$: Tensión eficaz máxima.  
    Máximo valor RMS de tensión inversa aplicable al diodo. 
    \item $P_{tot}$: Potencia total disipada.
    Máxima potencia que el dispositivo puede disipar sin dañarse.  
    \item $T_{rr}$: Tiempo de recuperación inversa.  
    Tiempo necesario para que el diodo deje de conducir al pasar de directa a inversa.  
\end{itemize}
